\chapter{Concluzii și direcții viitoare}

\section{Sinteza lucrării}
Lucrarea a avut ca obiectiv proiectarea și implementarea unui sistem de eLearning bazat pe blockchain, în limba română, conform cerințelor academice pentru o disertație. Tema a fost abordată prin analiza domeniului eLearning, a gamificării, a modelului learn-to-earn, a riscurilor recompensării prin tokenuri și a aplicațiilor hibride Web2/Web3, apoi prin proiectarea și documentarea platformei RustLearn. Rezultatul este un sistem care combină un backend Rust/Actix Web, o bază PostgreSQL, stocare S3, worker asincron, contracte Solidity și frontend Next.js.

Concluzia principală este că blockchain-ul aduce valoare într-o platformă educațională atunci când este folosit selectiv. Nu toate datele educaționale trebuie puse în blockchain. Lecțiile, progresul, permisiunile, evaluările, rapoartele și auditul operațional sunt mai potrivite într-o infrastructură controlată, rapidă și relațională. În schimb, tokenurile, transferurile și evenimentele financiare verificabile pot beneficia de transparența și interoperabilitatea Ethereum. Această separare este esențială pentru caracterul practic al soluției.

\section{Contribuții realizate}
Prima contribuție este definirea unei arhitecturi hibride pentru eLearning recompensat. Arhitectura arată cum pot coexista un sistem centralizat de cursuri și un strat blockchain de recompense fără ca unul să îl compromită pe celălalt. Backend-ul rămâne sursa de adevăr pentru reguli, iar blockchain-ul devine strat de verificare și transfer pentru tokenuri.

A doua contribuție este modelarea recompensei ca proces auditat. În loc să se acorde automat tokenuri la primul eveniment, RustLearn folosește candidați, politici, aprobări, planificare de plată, confirmare token, creditare portofel, notificare și reconciliere. Acest model este mai potrivit pentru o platformă reală deoarece acceptă erori, întârzieri și decizii administrative fără a pierde trasabilitatea.

A treia contribuție este separarea permisiunilor pe scopuri: platformă, organizație și curs. Această separare reflectă mai bine realitatea educațională decât un model global simplu. Un utilizator poate avea rol de profesor într-un curs, rol de membru într-o organizație și rol de cursant în alt context. Arhitectura permite exprimarea acestor situații fără privilegii excesive.

A patra contribuție este integrarea operațională a componentelor. Sistemul include migrații, verificări de sănătate, pregătire operațională, semnal periodic pentru worker, Docker Compose, teste și documentație de operare. Aceste elemente fac diferența între un prototip izolat și o platformă care poate fi rulată, verificată și extinsă.

\section{Răspuns la problema inițială}
Problema inițială a fost construirea unui sistem de eLearning bazat pe blockchain care să fie util, sigur și justificat arhitectural. Soluția propusă răspunde acestei probleme printr-o arhitectură modulară monolitică și hibridă. Monolitul modular reduce complexitatea operațională, iar separarea internă păstrează domeniile clare. Integrarea blockchain este limitată la tokenuri și fluxuri de valoare, ceea ce evită costuri inutile și protejează datele educaționale.

Tehnologiile alese sunt coerente cu această problemă. Rust oferă siguranță și performanță, Actix Web oferă un API modular, PostgreSQL și Diesel oferă consistență relațională, S3 susține conținut multimedia, worker-ul separă operațiile costisitoare, Next.js oferă suprafețe de produs clare, iar Solidity/OpenZeppelin oferă contracte mature pentru tokenuri. Fiecare tehnologie are un rol justificat în ansamblu.

\section{Limitări}
Lucrarea are limite care trebuie recunoscute. În primul rând, nu include un studiu cu utilizatori reali pentru a măsura impactul recompenselor asupra motivației. În al doilea rând, unele fluxuri sunt încă pregătite pentru dezvoltare viitoare, în special e-mailul real, încărcarea media avansată și unele detalii de progres persistent. În al treilea rând, contractele inteligente ar necesita audit extern înainte de utilizare în producție cu valoare reală. În al patrulea rând, orice sistem learn-to-earn trebuie analizat și din perspectivă juridică, fiscală și etică.

Aceste limitări nu invalidează soluția, ci indică diferența dintre o platformă aplicativă de disertație și un produs complet lansat public. Arhitectura a fost concepută tocmai pentru a permite aceste extinderi fără schimbări fundamentale.

\section{Direcții viitoare}
O direcție viitoare este introducerea unui serviciu real de e-mail, cu livrare, șabloane, reîncercări și monitorizare. O altă direcție este extinderea progresului persistent, astfel încât fiecare lecție, activitate și evaluare să contribuie la un model mai bogat de învățare. De asemenea, platforma poate adăuga analiză educațională pentru profesori și organizații, cu indicatori privind retenția, dificultatea lecțiilor și rata de promovare.

Pe partea blockchain, direcțiile viitoare includ audit extern pentru contracte, suport pentru rețele de test publice, mecanisme de trezorerie mai avansate, politici de emitere mai stricte și interfețe mai clare pentru MetaMask sau alte portofele. Reconcilierea poate fi extinsă cu sarcini programate, alerte și tablouri de bord pentru operatori. În plus, arderile de tokenuri și taxele de depunere sau retragere pot fi integrate într-un model economic mai formal.

Pe partea de produs, este necesar un studiu cu utilizatori reali. Platforma poate fi testată pe un set de cursuri pilot, comparând grupe cu recompense prin tokenuri și grupe fără recompense. Indicatorii relevanți ar putea include rata de finalizare, numărul de reveniri, rezultatele evaluărilor și percepția cursanților asupra utilității recompenselor. Acest studiu ar transforma lucrarea dintr-o contribuție tehnică într-o cercetare aplicată asupra impactului motivațional.

\section{Concluzie finală}
Sistemul de eLearning bazat pe blockchain prezentat în lucrare demonstrează că tokenizarea poate fi integrată într-o platformă educațională fără a sacrifica cerințele clasice de securitate, performanță și control. Valoarea soluției nu stă în folosirea blockchain-ului ca scop în sine, ci în folosirea lui acolo unde aduce transparență și verificabilitate. Prin arhitectura hibridă, RustLearn oferă o bază solidă pentru cursuri, evaluări, organizații, portofele și recompense LearnToken, păstrând posibilitatea de extindere către scenarii reale de învățare recompensată și cercetare aplicată asupra motivației cursanților.
